\usepackage[T1]{fontenc}
\usepackage{newtxtext}
\usepackage{amsmath,amssymb,mathtools,amsthm}
\usepackage{newtxmath}
\input{glyphtounicode}
\input{glyphtounicode-cmex}
\pdfglyphtounicode{barex}{007C}
\pdfglyphtounicode{parenlefttpA}{239B}
\pdfglyphtounicode{parenrighttpA}{239E}
\pdfglyphtounicode{parenleftbtA}{239D}
\pdfglyphtounicode{parenrightbtA}{23A0}
\pdfglyphtounicode{parenleftexA}{239C}
\pdfglyphtounicode{parenrightexA}{239F}
\pdfgentounicode=1
\usepackage[margin=1in]{geometry}
\usepackage{microtype}
\usepackage{booktabs,array}
\usepackage[dvipsnames]{xcolor}
\usepackage{graphicx,tikz}
\usepackage{enumitem,needspace}
\usepackage{textcomp,listings}
\usepackage[colorlinks=true,linkcolor=MidnightBlue,citecolor=MidnightBlue,urlcolor=MidnightBlue]{hyperref}
\hypersetup{pdftitle={Exact Fourier certificates for complex Hadamard matrices of order six},pdfauthor={OpenAI}}
\setlist{itemsep=2pt,topsep=5pt}
\newtheorem{theorem}{Theorem}[section]
\newtheorem{lemma}[theorem]{Lemma}
\newtheorem{proposition}[theorem]{Proposition}
\newtheorem{corollary}[theorem]{Corollary}
\theoremstyle{definition}
\newtheorem{definition}[theorem]{Definition}
\theoremstyle{remark}
\newtheorem{remark}[theorem]{Remark}
\newcommand{\C}{\mathbb C}
\newcommand{\T}{\mathbb T}
\newcommand{\Z}{\mathbb Z}
\newcommand{\Q}{\mathbb Q}
\newcommand{\E}{\mathbb E}
\DeclareMathOperator{\Ree}{Re}
\DeclareMathOperator{\ord}{ord}
\DeclareMathOperator{\tr}{Tr}
\DeclareMathOperator{\diag}{diag}
\setlength{\emergencystretch}{1em}
\lstdefinestyle{certificate}{
  language=Python,
  basicstyle=\fontsize{7.7}{9.1}\selectfont\ttfamily,
  keywordstyle=\bfseries,
  commentstyle=\color{black!65},
  stringstyle=\color{black},
  numbers=left,
  numberstyle=\fontsize{6.5}{8}\selectfont\color{black!60},
  numbersep=6pt,
  columns=fullflexible,
  keepspaces=true,
  showstringspaces=false,
  upquote=true,
  breaklines=true,
  breakatwhitespace=false,
  tabsize=4,
  xleftmargin=15pt,
  frame=none,
  aboveskip=10pt,
  belowskip=10pt
}
